% GENERATED by qa/build_cjk_reader_streams.py; DO NOT EDIT.
% Continuous preliminary matter with complete source milestones.
\VerdierSemanticAnchor{front-title}
\begin{titlepage}
  \VerdierSourcePageBlock{0002}
  \centering
  {\Huge\bfseries 239\hfill ASTÉRISQUE\par}
  \vspace{1.5em}
  {\LARGE\bfseries 1996\par}
  \vfill
  {\Huge\bfseries 阿贝尔范畴的\par}
  \vspace{0.4em}
  {\Huge\bfseries 导出范畴\par}
  \vspace{3em}
  {\Large\bfseries Jean-Louis VERDIER\par}
  \vfill
  {\Large\bfseries 法国数学会\par}
  \smallskip
  {\small\bfseries 本书由法国国家科学研究中心资助出版\par}
\end{titlepage}

\clearpage
\VerdierSemanticAnchor{front-ams-classification}
\VerdierSourcePageBlock{0003}
\VerdierReaderChapterStar{AMS 主题分类}
\addcontentsline{toc}{chapter}{AMS 主题分类}

14F20, 18E30, 18E35, 18G05, 18G10, 18G15, 18G20, 18G35,
18G40, 55U15, 55U25, 55U30.

\clearpage
\VerdierSemanticAnchor{front-portrait}
\VerdierSourcePageBlock{0004}
\thispagestyle{empty}
\vspace*{0.62\textheight}
\begin{center}
Jean-Louis VERDIER，1977
\end{center}

\VerdierSourcePageSilent{0005}%
\clearpage
\VerdierSemanticAnchor{front-preface}
\VerdierSourcePageBlock{0006}
\VerdierReaderChapterStar{序言}
\addcontentsline{toc}{chapter}{序言}

Jean-Louis Verdier 于 1989 年 8 月 25 日不幸逝世。他曾于 1967 年 6 月
14 日在巴黎答辩其国家博士论文，题为《阿贝尔范畴的导出范畴》。本书采用的
文本就是该论文打字稿的最终版本。编者所做的工作仅限于改正细小错误。提交给
学院的只有题为“引言”的部分；正文在答辩时尚未完成，此后也一直如此。

这项工作的由来已相当为人所知（\emph{参见} [I]）。我们知道，Grothendieck
在二十世纪六十年代初提出了导出范畴的思想，目的是为同调代数提供一个框架，
以实现他在 1958 年爱丁堡国际数学家大会报告之后所设想的对偶定理之宏大推广。
他把建立这一形式体系作为 Verdier 的博士论文课题。1963 年，Verdier 发表了
一册概述该理论核心内容的结果汇编《导出范畴：若干结果（状态 0）》，由 IHÉS
油印（后收入 [SGA 4 1/2]）。导出范畴的概念固然源自 Grothendieck，但更早
提出三角范畴概念的则是 Verdier。那里给出的公理体系受稳定同伦论启发
（\emph{参见} [P]），后来证明具有出人意料的丰硕成果。

又过了若干年，Verdier 才着手全面整理这项工作。他为什么没有完成？人们提出
过许多猜想。我倾向于以下解释：缺失的部分正是导出函子理论。当 Verdier 写到
这里时，导出范畴已经在代数几何中得到广泛使用，各种变体和推广也在发展：
Deligne 式的逐点可导性、非加性的导出函子（Quillen 式同伦代数），以及滤过
导出函子。此时，那册结果汇编所确定的框架已显得不足。要找到恰当的观点，必须
重新开展一项基础性工作；而当时正关注其他问题的 Verdier，在我看来既无暇，
大概也无意去完成它。随着时间推移，关于导出范畴的论文不断出现，其应用日益
普遍，这一形式体系还进入了数学的其他领域（例如线性偏微分方程）。于是，在
Verdier 眼中，完成这部著作越来越成为一项没有必要的工作……

\VerdierSourcePageBlock{0007}
这部文本虽不完整，却仍是一部珍贵乃至无可替代的参考资料（例如，其中对谱序列
所作的原创而优雅的论述，据我所知，在别处无从见到）。它终于得以问世，实在
令人欣慰。

\begin{flushright}
巴黎，1996 年 4 月 27 日\\
Luc Illusie
\end{flushright}

\section*{参考文献}

\begin{description}[labelwidth=6.2em,leftmargin=6.8em]
\item[{[I]}] L. Illusie，\emph{Catégories dérivées et dualité, travaux de
J.-L. Verdier}，L'Enseignement Mathématique，36 (1990)，369--391。
（1989 年 10 月 19 日，Université Paris VII 举办 Jean-Louis Verdier
纪念仪式时所作的报告。）

\item[{[P]}] D. Puppe，\emph{On the formal structure of stable homotopy
theory}，Colloquium on Algebraic Topology，Aarhus Universitet (1962)，
65--71。

\item[{\mbox{[SGA 4 1/2]}}] “Cohomologie étale”，Séminaire de Géométrie
Algébrique du Bois-Marie，SGA 4 1/2，P. Deligne 著，Lecture Notes in
Mathematics 569，Springer-Verlag (1977)，262--311。
\end{description}
\clearpage
\VerdierSemanticAnchor{front-editors-note}
\VerdierSourcePageBlock{0008}
\VerdierReaderChapterStar{编者按}
\addcontentsline{toc}{chapter}{编者按}

\VerdierSemanticAnchor{front-editors-note-publication-history}
1989 年 10 月 19 日，Université Paris VII 举行纪念 J.-L. Verdier 的仪式；
M. Artin、P. Cartier 和 L. Illusie 正是在那里萌生了出版 J.-L. Verdier
博士论文的想法。两年后，我加入他们，并在 J. Le Potier 与 M. Broué 的支持
下，从 Paris VII（URA 212）和 ÉNS（DMI）获得经费，将本文录入 \TeX{}。
录入完成后，我们把结果与 Verdier 的打字原稿作了仔细核对。由于 Verdier
从未作过最后修订，仍须进行一次数学内容上的审读。

\VerdierSemanticAnchor{front-editors-note-chapter-i-modifications}
这项细致的审读工作由 P. Cartier（第一章）、B. Keller（第三章）以及
M. Zisman 和我（全文）极其认真地完成。除大量不值得逐一列举的排印错误外，
主要改动如下。第一章命题 (3.1.5) 原来断言一个等价关系。我们把等价改为蕴含，
因为其逆命题（后文从未使用）是错误的。命题 (3.4.1) 有一个符号错误，并由此
造成第三章 ((1.3.4)、(b) 及 (3.2.7.3)) 的若干符号错误。这又要求调整“对偶
三角范畴”（第二章，1.1.7）的定义，以使下述断言仍然成立：\emph{对偶范畴的
导出范畴，就是原导出范畴的对偶三角范畴}。

\VerdierSemanticAnchor{front-editors-note-chapter-ii-modifications}
第二章中，§1.2 各命题的次序有所调整。Verdier 写完正文后认识到，两个特异
三角的直和仍为特异三角这一断言，可由三角范畴的其他公理推出。因此他把它从
公理表中删去，却没有完成由此产生的重新编排。此外，后文未使用且验证失败的
(2.3.6) 和 (2.3.7) 两段也被删除。定常谱对象 (4.4.2) 的定义作了轻微修改，
因为原定义过弱，不足以证明随后的定理 (4.4.3)。我们核实了这一修改不会影响
后续内容。

\VerdierSourcePageBlock{0009}
\VerdierSemanticAnchor{front-editors-note-apparatus}
记号索引和术语索引都得到了补充。与此相反，参考文献保持不变（但 SGA 各卷
改列其在 ``Lecture Notes'' 中的最终书目信息）。尤其是，参考文献编号 13 的
缺失并未补上。

\VerdierSemanticAnchor{front-editors-note-acknowledgements}
谨向 M. Broué、A. Bruguières、P. Cartier、L. Gruson、L. Illusie、
B. Keller、J. Le Potier、M. Zisman，以及 Université Paris VII、École
Normale Supérieure 和 Institut des Hautes Études Scientifiques 致以
诚挚谢意；他们以不同方式帮助本项目得以完成。我还感谢 C. Defosse 和
C. Gourgues 耐心地将本书录入 \textsc{AMS}-\TeX{}。

\VerdierSemanticAnchor{front-editors-note-memorial}
我本也希望能向 L. Doustaing 表达同样的谢意；遗憾的是，他已于 1994 年
12 月英年早逝。

\VerdierSemanticAnchor{front-editors-note-signature}
\begin{flushright}
巴黎，1996 年 6 月 14 日\\
Georges Maltsiniotis
\end{flushright}
\clearpage
\VerdierSemanticAnchor{front-publications}
\VerdierSourcePageBlock{0010}
\VerdierReaderChapterStar{J.-L. Verdier 出版物目录}
\addcontentsline{toc}{chapter}{J.-L. Verdier 出版物目录}

\small
\begin{description}
\VerdierSemanticAnchor{front-publications-01}
\item[{[01]}] \emph{与自守形式相联系的积分}（据 G. Shimura），Séminaire
Bourbaki 1960--61，no. 216，pp. 216-1--216-27，Benjamin (1966)。

\VerdierSemanticAnchor{front-publications-02}
\item[{[02]}] \emph{Poincaré 对偶定理}，C.R.A.S. vol. 256，
pp. 2084--2086，Gauthier-Villars (1963)。

\VerdierSemanticAnchor{front-publications-03}
\item[{[03]}] \emph{数域的 étale 上同调讨论班}（与 M. Artin 合作），
Woods Hole Conference on Algebraic Geometry，A.M.S. (1964)。

\VerdierSemanticAnchor{front-publications-04}
\item[{[04]}] \emph{投射系的本质等价}，C.R.A.S. vol. 261，
pp. 4950--4953，Gauthier-Villars (1965)。

\VerdierSemanticAnchor{front-publications-05}
\item[{[05]}] \emph{局部紧空间的维数}，C.R.A.S. vol. 261，
pp. 5293--5296，Gauthier-Villars (1965)。

\VerdierSemanticAnchor{front-publications-06}
\item[{[06]}] \emph{Profinite 群上同调中的对偶性}，收入 Jean-Pierre Serre
《Galois Cohomology》附录，Lecture Notes in Mathematics 5，
pp. 183--206，Springer (1964)，第四版 (1973)。

\VerdierSemanticAnchor{front-publications-07}
\item[{[07]}] \emph{局部紧空间上的可构造层}，C.R.A.S. vol. 262，
pp. 12--15，Gauthier-Villars (1966)。

\VerdierSemanticAnchor{front-publications-08}
\item[{[08]}] \emph{局部紧空间上同调中的对偶性}，Séminaire Bourbaki
1965--66，no. 300，pp. 300-01--300-13，Benjamin (1966)。

\VerdierSemanticAnchor{front-publications-09}
\item[{[09]}] \emph{étale 上同调中的 Lefschetz 不动点公式}，Conference
on Local Fields，1966 年 Driebergen NUFFIC Summer School，pp. 199--214，
Springer (1967)。

\VerdierSemanticAnchor{front-publications-10}
\item[{[10]}] \emph{概形的 étale 上同调中的一个对偶定理}，Conference
on Local Fields，1966 年 Driebergen NUFFIC Summer School，pp. 184--198，
Springer (1967)。

\VerdierSemanticAnchor{front-publications-11}
\item[{[11]}] \emph{凝聚层之扭曲逆像的基变换}，Bombay Colloquium on
Algebraic Geometry in 1968，pp. 393--408，Oxford University Press
(1969)。
\end{description}
\normalsize
\VerdierSourcePageContinue{0011}
\small
\begin{description}
\VerdierSemanticAnchor{front-publications-12}
\item[{[12]}] \emph{局部紧空间上同调的对偶定理}，Poincaré Duality，
Heidelberg-Strasbourg Seminar 1966--67，exp. 4，Publ. I.R.M.A.
Strasbourg，no. 3，pp. 1--36 (1969)。

\VerdierSemanticAnchor{front-publications-13}
\item[{[13]}] \emph{核定理}，Poincaré Duality，Heidelberg-Strasbourg
Seminar 1966--67，exp. 9，Publ. I.R.M.A. Strasbourg，no. 3，
pp. 1--52 (1969)。

\VerdierSemanticAnchor{front-publications-14}
\item[{[14]}] \emph{论 Wilder 的一个定理}，Applications of categorical
algebra，Proc. Symp. Pure Math. 17，pp. 184--191，A.M.S. (1970)。

\VerdierSemanticAnchor{front-publications-15}
\item[{[15]}] \emph{Grauert 凝聚性定理的一个简单证明}（与 R. Kiehl
合作），Math. Ann. 195，pp. 24--50，Springer (1971)。

\VerdierSemanticAnchor{front-publications-16}
\item[{[16]}] \emph{具有凝聚上同调的解析层复形之上同调空间上的拓扑}，
Bull. Soc. Math. France 99，pp. 337--343，Gauthier-Villars (1971)。

\VerdierSemanticAnchor{front-publications-17}
\item[{[17]}] \emph{复解析几何中的相对对偶}（与 J.-P. Ramis、G. Ruget
合作），Inv. Math. 13，pp. 261--283，Springer (1971)。

\VerdierSemanticAnchor{front-publications-18}
\item[{[18]}] S.G.A. 4 (1963--64)《拓扑斯理论与概形的 étale 上同调》：
\emph{预层}（与 A. Grothendieck 合作），pp. 1--184；\emph{拓扑与层}，
pp. 219--263；\emph{层范畴的函子性}，pp. 265--297；\emph{拓扑斯}
（与 A. Grothendieck 合作），pp. 299--519，Lecture Notes in Mathematics
269，Springer (1972)。

\VerdierSemanticAnchor{front-publications-19}
\item[{[19]}] S.G.A. 4 (1963--64)《拓扑斯理论与概形的 étale 上同调》：
\emph{拓扑斯中的上同调}，pp. 1--62；\emph{有限性条件、纤维化拓扑斯与
site 及其在取极限问题中的应用}（与 A. Grothendieck 合作），
pp. 163--336，Lecture Notes in Mathematics 270，Springer (1972)。

\VerdierSemanticAnchor{front-publications-20}
\item[{[20]}] \emph{Euler-Poincaré 示性数}，Bull. Soc. Math. France 101，
pp. 441--445，Gauthier-Villars (1973)。

\VerdierSemanticAnchor{front-publications-21}
\item[{[21]}] \emph{Le Potier 定理}，Seminar on Analytic Geometry
1972--73，Various aspects of positivity，Astérisque 17，pp. 68--78，
S.M.F. (1974)。

\VerdierSemanticAnchor{front-publications-22}
\item[{[22]}] \emph{$l$-进上同调中 Frobenius 自同态的特征多项式与 $l$
无关}（据 P. Deligne），Séminaire Bourbaki 1972--73，no. 423，
Lecture Notes in Mathematics 383，pp. 98--115，Springer (1974)。

\VerdierSemanticAnchor{front-publications-23}
\item[{[23]}] \emph{黎曼几何中的若干相交问题}（与 P. Marry 合作），
J. Diff. Geom. 12，pp. 345--376，A.M.S. (1977)。
\end{description}
\normalsize
\VerdierSourcePageContinue{0012}
\small
\begin{description}
\VerdierSemanticAnchor{front-publications-24}
\item[{[24]}] \emph{Whitney 分层与 Bertini-Sard 定理}，Inv. Math. 36，
pp. 295--312，Springer (1976)。

\VerdierSemanticAnchor{front-publications-25}
\item[{[25]}] \emph{与一个循环相联系的同调类}，Seminar on Analytic
Geometry 1974--75，Astérisque 36--37，pp. 101--151，S.M.F. (1976)。

\VerdierSemanticAnchor{front-publications-26}
\item[{[26]}] \emph{完全交的 Riemann-Roch 定理}，Seminar on Analytic
Geometry 1974--1975，Astérisque 36--37，pp. 189--228，S.M.F. (1976)。

\VerdierSemanticAnchor{front-publications-27}
\item[{[27]}] \emph{可带奇点的代数簇之 Riemann-Roch 定理}（据
P. Baum、W. Fulton 与 R. MacPherson），Séminaire Bourbaki 1974--75，
no. 464，Lecture Notes in Mathematics 514，pp. 159--175，Springer
(1976)。

\VerdierSemanticAnchor{front-publications-28}
\item[{[28]}] \emph{导出范畴}（1963 年写成，1976 年出版），Seminar on
Algebraic Geometry 4 1/2，Étale Cohomology，Lecture Notes in Mathematics
569，pp. 262--311，Springer (1977)。

\VerdierSemanticAnchor{front-publications-29}
\item[{[29]}] \emph{代数微分方程}，Séminaire Bourbaki 1977--78，no. 512，
Lecture Notes in Mathematics 710，pp. 101--122，Springer (1979)。

\VerdierSemanticAnchor{front-publications-30}
\item[{[30]}] \emph{瞬子}，The Yang-Mills equations，É.N.S. Seminar
1977--78，Astérisque 71--72，pp. 105--134，S.M.F. (1980)。

\VerdierSemanticAnchor{front-publications-31}
\item[{[31]}] \emph{Chern 类的特化}，Euler-Poincaré characteristic，
É.N.S. Seminar 1978--79，Astérisque 82--83，pp. 149--159，S.M.F. (1981)。

\VerdierSemanticAnchor{front-publications-32}
\item[{[32]}] \emph{Lie 代数、Hamilton 系统与代数曲线}（据 M. Adler
与 P. Van Moerbeke），Séminaire Bourbaki 1980--81，no. 566，Lecture
Notes in Mathematics 901，pp. 85--94，Springer (1981)。

\VerdierSemanticAnchor{front-publications-33}
\item[{[33]}] \emph{仿射 Lie 代数的表示：在若干物理问题中的应用}
（据 E. Date、M. Jimbo、M. Kashiwara 与 T. Miwa），Séminaire Bourbaki
1981--82，no. 596，Astérisque 92--93，pp. 365--377，S.M.F. (1982)。

\VerdierSemanticAnchor{front-publications-34}
\item[{[34]}] \emph{矩阵 $S$ 的构造}（与 A. Treibich-Kohn 合作），
Mathematics and Physics，É.N.S. Seminar 1979--82，Progress in Mathematics
37，pp. 175--196，Birkhäuser (1983)。

\VerdierSemanticAnchor{front-publications-35}
\item[{[35]}] \emph{二维 $\sigma$-模型及从 $S^2$ 到 $S^{2n}$ 的调和映射}，
Group Theoretical Methods in Physics，Lecture Notes in Physics 180，
pp. 136--141，Springer (1983)。

\VerdierSemanticAnchor{front-publications-36}
\item[{[36]}] \emph{微局部几何}，Algebraic Geometry，Japan-France
Conference，Tokyo-Kyoto 1982，Lecture Notes in Mathematics 1016，
pp. 125--133，Springer (1983)。
\end{description}
\normalsize
\VerdierSourcePageContinue{0013}
\small
\begin{description}
\VerdierSemanticAnchor{front-publications-37}
\item[{[37]}] \emph{McKay 对应的几何构造}（与 G. Gonzalez-Sprinberg
合作），Annales Scientifiques É.N.S. vol. 16，pp. 409--449，
Gauthier-Villars (1983)。

\VerdierSemanticAnchor{front-publications-38}
\item[{[38]}] \emph{层的特化与温和单值性}，Analyse et topologie sur les
espaces singuliers (II--III)，Astérisque 101--102，pp. 332--364，
S.M.F. (1983)。

\VerdierSemanticAnchor{front-publications-39}
\item[{[39]}] \emph{几何 Fourier 变换}（与 J.-L. Brylinski、
B. Malgrange 合作），C.R.A.S. vol. 297，pp. 55--58，
Gauthier-Villars (1983)。

\VerdierSemanticAnchor{front-publications-40}
\item[{[40]}] \emph{闭子空间上 pervers 层的延拓}，Systèmes Différentiels
et Singularités，Astérisque 130，pp. 210--217，S.M.F. (1985)。

\VerdierSemanticAnchor{front-publications-41}
\item[{[41]}] \emph{单值性 pervers 层的延拓}，Systèmes Différentiels et
Singularités，Astérisque 130，pp. 218--236，S.M.F. (1985)。

\VerdierSemanticAnchor{front-publications-42}
\item[{[42]}] \emph{有理二重点上的自反模}（与 M. Artin 合作），
Math. Ann. 270，pp. 79--82，Springer (1985)。

\VerdierSemanticAnchor{front-publications-43}
\item[{[43]}] \emph{正特征中的 McKay 规则}（与 G. Gonzalez-Sprinberg
合作），C.R.A.S. vol. 301，pp. 585--587，Gauthier-Villars (1985)。

\VerdierSemanticAnchor{front-publications-44}
\item[{[44]}] \emph{从 $S^2$ 到 $S^4$ 的调和映射}，Geometry today，
International Conference，Rome 1984，Progress in Mathematics 60，
pp. 267--282，Birkhäuser (1985)。

\VerdierSemanticAnchor{front-publications-45}
\item[{[45]}] \emph{有理二重点上自反模的乘法结构}（与
G. Gonzalez-Sprinberg 合作），Géométrie algébrique et applications：
I. Géométrie et calcul algébrique，Travaux en cours 22，pp. 79--110，
Hermann (1987)。

\VerdierSemanticAnchor{front-publications-46}
\item[{[46]}] \emph{一个令人伤感的计算}，附录，Géométrie algébrique et
applications：II. Singularités et géométrie complexe，Travaux en cours 23，
pp. 12--14，Hermann (1987)。

\VerdierSemanticAnchor{front-publications-47}
\item[{[47]}] \emph{几何 Fourier 变换 II}（与 J.-L. Brylinski、
B. Malgrange 合作），C.R.A.S. vol. 303，pp. 193--198，
Gauthier-Villars (1986)。

\VerdierSemanticAnchor{front-publications-48}
\item[{[48]}] \emph{量子群}（据 V. G. Drinfel'd），Séminaire Bourbaki
1986--87，no. 685，Astérisque 152--153，pp. 305--319，S.M.F. (1987)。

\VerdierSemanticAnchor{front-publications-49}
\item[{[49]}] \emph{从 $S^2$ 到 $S^4$ 的调和映射 II}，Harmonic mappings,
twistors, and $\sigma$-models，Advanced series in Mathematical Physics 4，
pp. 124--147，World Scientific (1988)。
\end{description}
\normalsize
\VerdierSourcePageContinue{0014}
\small
\begin{description}
\VerdierSemanticAnchor{front-publications-50}
\item[{[50]}] \emph{新椭圆孤子}，``Algebraic Analysis''，Papers Dedicated
to Professor Mikio Sato on the Occasion of his Sixtieth Birthday，Vol. II，
pp. 901--910，Academic Press (1988)。

\VerdierSemanticAnchor{front-publications-51}
\item[{[51]}] \emph{椭圆孤子}（与 A. Treibich 合作），The Grothendieck
Festschrift，Vol. III，Progress in Mathematics 88，pp. 437--473，
Birkhäuser (1990)。

\VerdierSemanticAnchor{front-publications-52}
\item[{[52]}] \emph{切向覆盖与四个三角数之和}（与 A. Treibich 合作），
C.R.A.S. vol. 311，pp. 51--54，Gauthier-Villars (1990)。

\VerdierSemanticAnchor{front-publications-53}
\item[{[53]}] \emph{Boussinesq 椭圆孤子：循环情形}（与 E. Previato
合作），Proceedings of the Indo-French Conference on Geometry，
pp. 173--185，Hindustan Book Agency (1993)。

\VerdierSemanticAnchor{front-publications-54}
\item[{[54]}] \emph{K P 椭圆孤子的 Kritchever 簇}（与 A. Treibich
合作），Proceedings of the Indo-French Conference on Geometry，
pp. 187--232，Hindustan Book Agency (1993)。
\end{description}
\normalsize
\VerdierSourcePageSilent{0015}%
\clearpage
\begingroup
  \thispagestyle{empty}
  \centering
  \VerdierSemanticAnchor{thesis-title}
  \VerdierSourcePageBlock{0016}
  \vspace*{0.24\textheight}
  {\Huge 阿贝尔范畴的\par}
  \vspace{1.6em}
  {\Huge 导出范畴\par}
  \vspace{7em}
  {\Large\bfseries J.-L. VERDIER\par}
  \vfill
\par
\endgroup
\VerdierSourcePageSilent{0017}%
